\subsection{Agents and their environments}
\label{sec:agents-and-environments}
\slidesrc{Agent architectures: 1--10}

\subsubsection{Three classic definitions of agent}
\slidesrc{Agent architectures: 3--5}

The slides recall three classic definitions (the sources are added):
\begin{enumerate}
  \item ``An agent is anything that can be viewed as \emph{perceiving its environment
        through sensors} and \emph{acting upon that environment through effectors}.''
        \hfill (Russell and Norvig)
  \item ``Autonomous agents are computational systems that inhabit some complex dynamic
        environment, sense and act autonomously in this environment, and by doing so realize
        a set of goals or tasks for which they are designed.'' \hfill (P.~Maes)
  \item ``An autonomous agent is a system situated within and a part of an environment that
        senses that environment and acts on it, over time, in pursuit of its own agenda and
        so as to affect what it senses in the future.'' \hfill (Franklin and Graesser)
\end{enumerate}
All three share the \textbf{basic abstract view} of slide 4: a loop between agent and
environment, through sensors (input) and effectors (output); the same picture as
Section~\ref{sec:ai-definition}. The third definition adds an important nuance:
the agent's actions change what it will \emph{perceive in the future}.

\begin{nbdef}{Reactivity (slide 5)}
An agent has to be able to \textbf{react}, i.e.\ adapt its behaviour, in an appropriate way
to the dynamic changes in its environment: other computational agents, human
agents/users, external information sources (e.g.\ sensors), physical objects (e.g.\
robots), the Internet, \ldots{} Reactivity is one of several properties an intelligent agent
should have (Section~\ref{sec:what-is-an-agent}).
\end{nbdef}

\subsubsection{Kinds of environments}
\slidesrc{Agent architectures: 6--10}

The difficulty of building an agent depends heavily on the environment it lives in. Five
dimensions classify environments; in each pair, the first is the \emph{easy} one.

\begin{center}
\small
\renewcommand{\arraystretch}{1.2}
\begin{tabularx}{\linewidth}{@{}>{\raggedright\arraybackslash}p{2.6cm} >{\raggedright\arraybackslash}X >{\raggedright\arraybackslash}X@{}}
\toprule
\textbf{Dimension} & \textbf{Easy case} & \textbf{Hard case} \\
\midrule
accessible vs.\ inaccessible & the agent can obtain \emph{complete, accurate, up-to-date}
  information about the environment's state & most moderately complex environments
  (the physical world, the Internet) are inaccessible \\
deterministic vs.\ non-deterministic & any action has a \emph{single guaranteed effect}: no
  uncertainty about the resulting state & the physical world is, for all intents and
  purposes, non-deterministic \\
episodic vs.\ non-episodic & performance depends on discrete \emph{episodes} with no link
  between them: the agent decides based only on the current episode & the agent must reason
  about how current actions affect future episodes \\
static vs.\ dynamic & remains unchanged except by the agent's own actions & other processes
  operate on it and change it beyond the agent's control; the physical world is highly
  dynamic \\
discrete vs.\ continuous & a fixed, finite number of actions and percepts & the real
  world is continuous \\
\bottomrule
\end{tabularx}
\captionof{table}{Kinds of environments (slides 6--10). The more an environment is on the
left, the simpler it is to build agents for it.}
\end{center}

\begin{nbexercise}{E.1}{Practical exercise, step 1: think about the kind of environment (slide 10)}
Classify these environments along the five dimensions:
(a) a chess program (with a clock); (b) an e-mail spam filter; (c) a self-driving taxi;
(d) a robot vacuum cleaner; (e) an online poker bot.
\end{nbexercise}
\begin{nbsolution}
\small
\begin{tabular}{@{}lccccc@{}}
\toprule
 & \textbf{access.} & \textbf{determ.} & \textbf{episodic} & \textbf{static} & \textbf{discrete} \\
\midrule
(a) chess with clock & yes & yes & no & semi$^\ast$ & yes \\
(b) spam filter & no & no & yes & no & yes \\
(c) self-driving taxi & no & no & no & no & no \\
(d) robot vacuum & no & no & no & no & no \\
(e) poker bot & no & no & no$^{\ast\ast}$ & no & yes \\
\bottomrule
\end{tabular}

\normalsize
Explanations: chess is fully observable and every move has a certain effect; the board
does not change while you think, but the clock does ($^\ast$semi-dynamic). A spam filter
classifies each e-mail independently (episodic), but it cannot see the sender's intentions
(inaccessible). The taxi is the hardest case on every dimension, like the vacuum robot (the
cat moves, the floor is continuous). Poker hides the opponents' cards (inaccessible), the
deck is shuffled (non-deterministic), and ($^{\ast\ast}$) hands are episodes, but the
opponents' model of you carries over from hand to hand. \textbf{Lesson}: almost every
real application is on the ``hard'' side, which is why architectures matter.
\end{nbsolution}
